\documentclass[letterpaper,10pt]{article}
\usepackage{latexsym}
\usepackage[empty]{fullpage}
\usepackage{titlesec}
\usepackage{marvosym}
\usepackage[usenames,dvipsnames]{color}
\usepackage{xcolor}
\usepackage{verbatim}
\usepackage{enumitem}
\usepackage[hidelinks]{hyperref}
\usepackage[english]{babel}
\usepackage{tabularx}
\usepackage{iftex}
\ifPDFTeX
  \input{glyphtounicode}
\fi

% Original resume typography
\usepackage[T1]{fontenc}
\usepackage{charter}

\definecolor{cvblue}{HTML}{0E5484}
\definecolor{black}{HTML}{130810}
\definecolor{darkcolor}{HTML}{0F4539}
\definecolor{cvgreen}{HTML}{3BD80D}
\definecolor{taggreen}{HTML}{00E278}
\definecolor{SlateGrey}{HTML}{2E2E2E}
\definecolor{LightGrey}{HTML}{666666}
\colorlet{name}{black}
\colorlet{tagline}{darkcolor}
\colorlet{heading}{darkcolor}
\colorlet{headingrule}{cvblue}
\colorlet{accent}{darkcolor}
\colorlet{emphasis}{SlateGrey}
\colorlet{body}{LightGrey}

% Adjust margins
\addtolength{\oddsidemargin}{-0.6in}
\addtolength{\evensidemargin}{-0.5in}
\addtolength{\textwidth}{1.19in}
\addtolength{\topmargin}{-.7in}
\addtolength{\textheight}{1.4in}
\urlstyle{same}

\raggedbottom
\raggedright
\setlength{\tabcolsep}{0in}
\ifPDFTeX
  \pdfgentounicode=1
\fi
\ifXeTeX
  \XeTeXgenerateactualtext=1
\fi

% Section formatting
\titleformat{\section}{
  \vspace{-2pt}\scshape\raggedright\large\bfseries
}{}{0em}{}[\color{black}\titlerule \vspace{-2pt}]

% Custom commands
\newcommand{\resumeItem}[1]{\item\small{#1 \vspace{-2pt}}}
\newcommand{\resumeSubheading}[4]{
  \vspace{-2pt}\item
    \begin{tabular*}{1.0\textwidth}[t]{l@{\extracolsep{\fill}}r}
      \textbf{\large #1} & {\small #2} \\
      \textit{\small #3} & \textit{\small #4} \\
    \end{tabular*}\vspace{-2pt}
}
\newcommand{\resumeSubItem}[1]{\resumeItem{#1}\vspace{0pt}}
\newcommand{\resumeItemListStart}{\begin{itemize}[leftmargin=0.15in, label={\color{cvblue}-}]}
\newcommand{\resumeItemListEnd}{\end{itemize}\vspace{0pt}}
\newcommand{\resumeSubHeadingListStart}{\begin{itemize}[leftmargin=0.0in, label={}]}
\newcommand{\resumeSubHeadingListEnd}{\end{itemize}}

%-------------------------------------------
%%%%%%  RESUME STARTS HERE  %%%%%%%%%%%%%%%%%%%%%%%%%%%%

\begin{document}

%----------HEADING----------
\begin{center}
    {\huge \textbf{Haixi Zhang}} \\
    \vspace{3pt}
    \href{https://www.linkedin.com/in/haixi-zhang-b83104251/}{\color{blue}{LinkedIn}}
    \small{-}
    \href{mailto:haixizhang02@gmail.com}{\color{blue}{haixizhang02@gmail.com}}
    \small{-} {\color{blue}{(607)\,339-9816}}
    \small{-} {\color{blue}{San Jose, CA 95113}}
    \small{-}
    \href{https://haixizhang.github.io/}{\color{blue}{Portfolio}}
    \small{-}
    \href{https://github.com/haixizhang}{\color{blue}{GitHub}}
\end{center}

%-----------PROFILE-----------
% \section{\color{cvblue}PROFILE}
% \small{
% Perception/Robotics engineer focused on \textbf{C++/Python}, building \textbf{camera–LiDAR} perception and calibration pipelines on \textbf{Linux/ROS}. Hands-on experience in \textbf{2D/3D detection, tracking, depth/point-cloud processing}, and \textbf{sensor fusion}, with deployment via Docker and system-level testing. Comfortable with real-time constraints and cross-team integration.
% }
% \section{\color{cvblue}PROFILE}
% \small{
% Sensor/ML engineer focused on \textbf{estimation and sensor-fusion algorithms} for robotics and embedded platforms.
% Strong in \textbf{linear algebra, probability/statistics}, and \textbf{camera/LiDAR geometry and calibration}.
% Hands-on experience building end-to-end pipelines from \textbf{data collection \& algorithm validation} (Python/PyTorch/OpenCV) to \textbf{C/C++ integration} on Linux/ROS, with attention to \textbf{latency and compute/memory tradeoffs}.
% }

%-----------EDUCATION-----------
\section{\color{cvblue}EDUCATION}
\resumeSubHeadingListStart
  \resumeSubheading
    {Cornell University}{Ithaca, NY}
    {M.Eng in Electrical \& Computer Engineering (GPA: 3.81/4.0)}{Aug. 2024 -- May. 2025}
  \resumeSubheading
    {University of Rochester}{Rochester, NY}
    {B.S. in Electrical \& Computer Engineering; Minor in Statistics (GPA: 3.86/4.0)}{Sep. 2020 -- May. 2024}
\resumeSubHeadingListEnd

%-----------TECHNICAL SKILLS-----------
% \section{\color{cvblue}TECHNICAL SKILLS}
% \small{
% \textbf{Languages:} C, Python, Java, C++, Bash, MATLAB, R \\
% \textbf{Embedded \& OS:} Embedded Linux, ROS, RTOS, ARM Cortex Microcontrollers, Sensors (ToF, IMU, LiDAR), Mixed-Signal Circuits, TCP/IP protocols,
% Serial Protocols, BLE, Hardware Integration, SIL/HIL Testing, Schematics/PCB Reading, Windows  \\
% \textbf{Tools:} CI/CD (GitLab), Git, Docker, Kubernetes, Logic Analyzers, Oscilloscope, Digital Multimeters \\
% \textbf{Data \& ML:} PyTorch, HPC, Computer Vision, NLP, Nonlinear Optimization, Statistical Analysis, Multiprocessing \\
% }

\section{\color{cvblue}TECHNICAL SKILLS}
\small{
\textbf{Languages \& Systems:} C++, Python, CUDA, Bash, Linux, ROS/ROS2, Git, Bazel, CMake, Docker, Kubernetes \\
\textbf{Computer Vision \& Robotics:} Camera models, multi-sensor extrinsic calibration, $SE(3)/SO(3)$, LiDAR--camera fusion, point-cloud processing, OpenCV, Eigen, PCL/Open3D, state estimation (EKF), SLAM, occupancy-grid mapping, A* planning \\
\textbf{Machine Learning \& AI:} PyTorch, CNNs, Vision Transformers, multimodal learning, cross-attention, multi-task learning, representation learning, PyTorch DistributedDataParallel (DDP), NCCL, distributed data loading/evaluation, LLM APIs, AI agents \\
\textbf{Production ML \& Data:} ONNX, TensorRT, C++ inference, model export and optimization, numerical-parity testing, Airflow, ETL pipelines, dataset curation, data-quality validation, production monitoring
}

%-----------PROFESSIONAL EXPERIENCE-----------
\section{\color{cvblue}PROFESSIONAL EXPERIENCE}
\resumeSubHeadingListStart
\resumeSubheading
  {Machine Learning Engineer, Tensor Auto}{San Jose, CA}
    {Production Computer Vision, 3D Geometry, Multimodal Learning, Onboard Inference}{Sep. 2025 -- Present}
    \resumeItemListStart
      \resumeItem{Own three company-level ML/CV systems as sole technical owner, taking two from initial requirements to onboard production deployment and leading the third from requirements through active validation and integration.}
      \resumeItem{Designed a residual extrinsic-calibration framework spanning camera--LiDAR 6-DoF alignment and LiDAR--LiDAR rotational alignment, combining differentiable projection and depth rasterization, $SE(3)/SO(3)$ correction modeling, multimodal feature fusion, and geometry-aware objectives.}
      \resumeItem{Co-developed Airflow/Kubernetes data pipelines and curated 50K+ training pairs from 500+ GB of multimodal sensor logs; both calibration systems recovered controlled rotational perturbations up to $10^\circ$ to approximately $0.5^\circ$ residual, while LiDAR--LiDAR evaluation improved rotation by 62.5\% across 95.9\% of held-out frames.}
      \resumeItem{Re-engineered CPU-heavy camera--LiDAR training via geometry precomputation and TAR-sharded streaming, reducing single-accelerator runs from $\sim$60 to $\sim$20 hours; built world-size-configurable PyTorch DistributedDataParallel/NCCL with deterministic sharding and globally correct evaluation to reach $<$5 hours ($\sim$13$\times$ overall), while profiling LiDAR--LiDAR scale-out limits to CPU-side projection and point correspondence.}
      \resumeItem{Lead development of an in-cabin seat-occupancy and occupant-posture classification system from fisheye imagery, establishing session-isolated evaluation, multi-seed benchmarking, class-imbalance diagnostics, deterministic output gating, and an ONNX-compatible production contract.}
    \resumeItemListEnd

  \resumeSubheading
    {Undergraduate Research Assistant}{University of Rochester, ECE Department}
    {Python, PyTorch, C, Graph Learning, HPC, Docker/Kubernetes}{Jan. 2023 -- May 2024}
    \resumeItemListStart
    \resumeItem{Built reproducible HPC pipelines for graph learning and nature-inspired models in financial and power-system applications, improving benchmark performance by up to 25\% and containerizing distributed experiments with Docker/Kubernetes.}
    \resumeItemListEnd
\resumeSubHeadingListEnd
%-----------PROJECTS \& EXPERIENCE-----------
\section{\color{cvblue}SELECTED RESEARCH \& PROJECTS}
\resumeSubHeadingListStart
  \resumeSubheading
    {Robotic Tool Handler (Senior Design)}{University of Rochester, Lead Developer}
    {C++, ROS, NVIDIA Jetson, Real-time Perception \href{https://github.com/haixizhang/Robotic-Tool-Handler}{\color{blue}\underline{GitHub}}}{Jan. 2024 -- Jun. 2024}
    \resumeItemListStart
      \resumeItem{Led perception and system integration for a seven-member robotics team; deployed a containerized ROS vision pipeline on NVIDIA Jetson, improved throughput by 20\%, and built GitLab CI with software- and hardware-in-the-loop validation.}
      \resumeItem{Developed real-time 3D object detection and classification by fusing camera and LiDAR evidence, achieving sub-20 ms inference latency and improving mean average precision by 25\%.}
    \resumeItemListEnd

  \resumeSubheading
    {Robot Mapping, Estimation, and Navigation}{University of Rochester, Solo Developer}
    {ROS, C++, LiDAR, Kalman Filter, TurtleBot2 \href{https://github.com/haixizhang/Robot-Mapping-Estimation-Interaction}{\color{blue}\underline{GitHub}}}{Jan. 2024 -- May. 2024}
    \resumeItemListStart
      \resumeItem{Implemented and deployed a C++/ROS navigation stack on TurtleBot2, integrating RGB-D, LiDAR, and wheel odometry for EKF-based state estimation, occupancy-grid mapping, A* planning, and sensor fault detection.}
      \resumeItem{Reduced localization error by 10\% and planning time by 15\% through EKF tuning, occupancy-map parameterization, and unit-tested sensor integration and failure handling.}
    \resumeItemListEnd

  \resumeSubheading
    {Monocular Real-time Perception for Autonomous Driving}{University of Rochester}
    {Python, PyTorch, Computer Vision, Depth Estimation \href{https://github.com/EnisZuo/LRHPerception}{\color{blue}\underline{GitHub}}}{Apr. 2023 -- Aug. 2023}
    \resumeItemListStart
      \resumeItem{Designed a 30+ FPS multi-task perception pipeline combining object tracking, trajectory prediction, road segmentation, and monocular depth estimation through shared visual features and lightweight task decoders.}
    \resumeItemListEnd
\resumeSubHeadingListEnd



\end{document}
